\documentclass[11pt]{article}
\usepackage[margin=1in]{geometry}
\usepackage[hidelinks]{hyperref}
\usepackage{parskip}
\usepackage[sfdefault]{roboto}
\pagenumbering{gobble}

\begin{document}

\textbf{Gur Raunaq Singh}\\
Bangalore, India\\
+91-9015154091\\
\href{mailto:raunaq.soni@gmail.com}{raunaq.soni@gmail.com}\\
\href{https://www.linkedin.com/in/gurraunaqsingh/}{LinkedIn} \textbar{} \href{https://github.com/raunaqness}{GitHub}

March 5, 2026

Hiring Team,

I am applying for the Senior Generative AI Engineer / AI Founding Engineer role at AISEO. I bring 8+ years of backend engineering experience and recent hands-on work shipping LLM-powered systems in production, including RAG pipelines, prompt evaluation workflows, and cloud-native AI infrastructure.

At Aptos India, I built enterprise AI systems for high-volume support workflows. I designed RAG pipelines with vector retrieval and context selection, implemented prompt versioning with quality evaluation, and reduced hallucinations by 35\% through iterative testing and runtime monitoring. I also architected secure GCP data/model pipelines with IAM and audit controls, improving SLA performance by 35\% and reducing compute spend by 60\%.

On fine-tuning and model adaptation: my production work has primarily focused on prompt, retrieval, and evaluation optimization layers around foundation models, while applying model behavior tuning through controlled experimentation and domain-specific data curation. I am comfortable extending this into adapter-based workflows (for example LoRA/QLoRA) for SEO-specific tasks where measurable gains justify specialization.

A relevant technical challenge I solved was improving answer quality in enterprise knowledge-heavy workflows where factuality and usefulness were both critical. I addressed this by combining semantic retrieval tuning, prompt templates with strict grounding constraints, and offline/online evaluation loops. The result was better response quality with lower hallucination rates and stronger operational trust.

My GEO perspective is that search is shifting from link-ranking-first to answer-composition-first. Winning content must be machine-interpretable, entity-rich, and semantically aligned to intent across both traditional SERPs and AI answer engines. I would approach a humanization system by separating concerns: generation, authenticity scoring, SEO scoring, and policy checks, then optimizing each stage with measurable metrics (readability, factuality, topical coverage, and ranking performance) instead of a single black-box prompt.

I am available for full-time exclusive commitment and can align with Amsterdam working hours.

Portfolio links:

\begin{itemize}
  \item GitHub: \url{https://github.com/raunaqness}
  \item AI Travel Agent repository: \url{https://github.com/raunaqness/ai-travel-agent}
  \item Live demo: \url{https://ai-travel-agent.fly.dev/}
\end{itemize}

I would welcome the opportunity to help AISEO build category-defining SEO/GEO products for the AI-first search era.

Sincerely,\\
Gur Raunaq Singh

\end{document}
